% TODO make hw stylesheet
\documentclass[10pt]{article}
\usepackage[letterpaper,top=1in,left=1in,right=1in]{geometry}
\usepackage[dvipsnames,svgnames]{xcolor}
\usepackage[shortlabels]{enumitem}
\usepackage[framemethod=TikZ]{mdframed}
\usepackage{amsmath,amssymb,amsthm}
\usepackage[colorlinks]{hyperref}
\usepackage{multicol}
\usepackage{tabularx}
\usepackage{bbm}

\hypersetup{
    colorlinks=true,
    linkcolor=blue,
    filecolor=magenta,
    urlcolor=magenta,
    pdftitle={MAT 549 HW}
}

\usepackage{mathtools}
\usepackage{thmtools}
\usepackage[textsize=footnotesize]{todonotes}
\usepackage{titling}
\usepackage{cancel}

\reversemarginpar
\mdfdefinestyle{mdbluebox}{roundcorner=10pt,innerbottommargin=9pt,
linecolor=blue,backgroundcolor=TealBlue!5,}
\declaretheoremstyle[headfont=\sffamily\bfseries\color{MidnightBlue},
mdframed={style=mdbluebox},]{thmbluebox}

\mdfdefinestyle{mdbloobox}{frametitlefont=\bfseries,innerbottommargin=8pt,
nobreak=true,backgroundcolor=TealBlue!5,linecolor=blue,}
\declaretheoremstyle[headfont=\bfseries\color{MidnightBlue},
mdframed={style=mdbloobox},headpunct={\\[3pt]},postheadspace=0pt,]{thmbloobox}

\mdfdefinestyle{mdredbox}{frametitlefont=\bfseries,innerbottommargin=8pt,
nobreak=true,backgroundcolor=Salmon!5,linecolor=RawSienna,}
\declaretheoremstyle[headfont=\bfseries\color{RawSienna},
mdframed={style=mdredbox},headpunct={\\[3pt]},postheadspace=0pt,]{thmredbox}

\mdfdefinestyle{mdgreenbox}{linecolor=ForestGreen,backgroundcolor=ForestGreen!5,
linewidth=2pt,rightline=false,leftline=true,topline=false,bottomline=false,}
\declaretheoremstyle[headfont=\bfseries\sffamily\color{ForestGreen!70!black},
mdframed={style=mdgreenbox},headpunct={ --- },]{thmgreenbox}

\mdfdefinestyle{mdorangebox}{linecolor=RedOrange,backgroundcolor=RedOrange!5,
linewidth=2pt,rightline=false,leftline=true,topline=false,bottomline=false,}
\declaretheoremstyle[headfont=\bfseries\sffamily\color{RedOrange!70!black},
mdframed={style=mdorangebox},headpunct={ --- },]{thmorangebox}

\mdfdefinestyle{mdblackbox}{linecolor=black,backgroundcolor=RedViolet!5!gray!5,
linewidth=3pt,rightline=false,leftline=true,topline=false,bottomline=false,}
\declaretheoremstyle[mdframed={style=mdblackbox}]{thmblackbox}

\declaretheoremstyle[spaceabove=6pt,spacebelow=6pt]{boldhead}
\declaretheoremstyle[spaceabove=6pt,spacebelow=6pt,headfont=\normalfont\itshape]{italichead}

\declaretheorem[style=boldhead,name=Problem]{problem}
\declaretheorem[style=boldhead,name=Problem,numbered=no]{problem*}
\declaretheorem[style=boldhead,name=Setup]{setup} % for config geo
\declaretheorem[style=boldhead,name=Example,sibling=problem]{example}
\declaretheorem[style=boldhead,name=$\bigstar$ Problem,sibling=problem]{niceprob}
\declaretheorem[style=boldhead,name=Theorem,sibling=problem]{theorem}
\declaretheorem[style=boldhead,name=Corollary,sibling=theorem]{corollary}
\declaretheorem[style=boldhead,name=Proposition,sibling=theorem]{prop}
\declaretheorem[style=boldhead,name=Lemma,numberwithin=problem]{lemma}
\declaretheorem[style=boldhead,name=Theorem,numbered=no]{theorem*}
\declaretheorem[style=boldhead,name=Lemma,numbered=no]{lemma*}
\declaretheorem[style=boldhead,name=Claim,numberwithin=problem]{claim}
\declaretheorem[style=boldhead,name=Claim,numbered=no]{claim*}
\declaretheorem[style=boldhead,name=Remark,numberwithin=problem]{remark}
\declaretheorem[style=boldhead,name=Remark,numbered=no]{remark*}
\declaretheorem[style=boldhead,name=Definition,sibling=theorem]{definition}
\declaretheorem[style=boldhead,name=Definition,numbered=no]{definition*}

\providecommand{\ol}{\overline}
\providecommand{\ceil}[1]{\left\lceil #1 \right\rceil}
\providecommand{\floor}[1]{\left\lfloor #1 \right\rfloor}
\newcommand{\dd}{\mathrm{d}}
\providecommand{\ul}{\underline}
\providecommand{\eps}{\varepsilon}
\providecommand{\half}{\frac{1}{2}}
\providecommand{\CC}{\mathbb C}
\providecommand{\FF}{\mathbb F}
\providecommand{\II}{\mathbb I}
\providecommand{\NN}{\mathbb N}
\providecommand{\QQ}{\mathbb Q}
\providecommand{\RR}{\mathbb R}
\providecommand{\ZZ}{\mathbb Z}
\providecommand{\ind}{\mathbbm 1}
\providecommand{\dg}{^\circ}
\providecommand{\ii}{\item}
\providecommand{\setdiff}{\smallsetminus}
\providecommand{\alert}[1]{{\sffamily\textbf{\textcolor{blue}{#1}}}}
\providecommand{\inv}{^{-1}}
\providecommand{\mb}[1]{\mathbf{#1}}
\newcommand{\what}{\widehat}

\newenvironment{soln}{\begin{proof}[Solution]}{\end{proof}}

\providecommand{\pow}{\operatorname{Pow}}
\providecommand{\vocab}[1]{{\textbf{\textcolor{ForestGreen}{#1}}}}
\providecommand{\lcm}{\operatorname{lcm}}
\providecommand{\abs}[1]{\left\lvert #1\right\rvert}
\providecommand{\smabs}[1]{\lvert #1\rvert}
\providecommand{\innerprod}[1]{\langle\!\langle #1\rangle\!\rangle}
\providecommand{\norm}[1]{\left\| #1\right\|}
\providecommand{\smnorm}[1]{\| #1\|}
\providecommand{\gr}{\operatorname{gr}}

\usepackage{mathrsfs}
\providecommand{\tanvec}[1]{\frac{\partial}{\partial #1}}

% place title at top right
\setlength{\droptitle}{-8ex}
\pretitle{\begin{flushright}\Large\bfseries}
\posttitle{\par\end{flushright}}
\preauthor{\begin{flushright}}
\postauthor{\end{flushright}}
\predate{\begin{flushright}}
\postdate{\end{flushright}}

\title{MAT 549 HW07}
\author{\large Aditya Pahuja}
\date{\large Due 04/28}
\begin{document}
\maketitle

\begin{problem}
    [17-1]
    Let $M$ be a smooth manifold with or without boundary
    and let $\omega\in\Omega^p(M)$, $\eta\in\Omega^q(M)$
    be closed forms. Show that the de Rham cohomology class
    of $\omega\wedge\eta$ depends only on
    the cohomology classes of $\omega$ and $\eta$,
    and thus there is a well-defined bilinear map
    $\smile\colon H^p_{dR}(M)\times H^q_{dR}(M)\to H^{p+q}_{dR}(M)$,
    called the cup product, given by
    $[\omega]\smile[\eta] = [\omega\wedge\eta]$.
\end{problem}

\begin{soln}
    Let $\omega'\in[\omega]$ and $\eta'\in[\eta]$,
    so $\omega' - \omega = \dd\alpha$ and $\eta' - \eta = \dd\beta$ are exact.
    Then
    \begin{equation*}
	\omega\wedge\eta - \omega'\wedge\eta'
	= (\omega - \omega')\wedge\eta
	+ \omega'\wedge(\eta - \eta')
	= \dd\alpha\wedge\eta + \omega'\wedge\dd\beta
	= \dd(\alpha\wedge\eta + \omega'\wedge\beta),
    \end{equation*}
    where the last equality uses the fact that
    $\eta$ and $\omega'$ are closed, so
    $[\omega\wedge\eta] = [\omega'\wedge\eta']$.
    Thus $\smile$ is a well-defined map on the cohomology groups,
    and it inherits bilinearity from the bilinearity of $\wedge$:
    explicitly, we can write
    \begin{equation*}
	[a\omega + b\eta]\smile[\alpha] = [(a\omega + b\eta)\wedge\alpha]
	= [a\omega\wedge\alpha] + [b\eta\wedge\alpha]
	= a[\omega]\smile[\alpha] + b[\eta]\smile[\alpha]
    \end{equation*}
    to show linearity in the first argument,
    and essentially the same calculation
    shows linearity in the second argument.
\end{soln}

\begin{problem}
    [17-12]
    Suppose $M$ and $N$ are compact, connected, and oriented $n$-manifolds,
    and $F\colon M\to N$ is a smooth map.
    Prove that if $\int_M F^*\eta\ne 0$ for some $\eta\in\Omega^n(N)$,
    then $F$ is surjective. Give an example to show that
    $F$ can be surjective even if $\int_M F^*\eta = 0$
    for every $\eta\in\Omega^n(N)$.
\end{problem}

\begin{soln}
    Suppose $F$ is not surjective.
    Then, any point not in $F(M)$ is vacuously a regular value for $F$,
    so the degree of $F$ is zero; that is,
    \begin{equation*}
        \int_M F^*\eta = \deg F \cdot \int_N \eta = 0
    \end{equation*}
    for all $\eta\in\Omega^n(N)$.

    To show that the converse does not hold,
    let $M = N = S^1\cong \RR/\ZZ$ and consider
    $F\colon S^1\to S^1$ given by
    \begin{equation*}
	F(x) = \sin^2(\pi x).
    \end{equation*}
    Since $x\mapsto\sin^2(\pi x)$ is a surjective smooth map $\RR\to[0,1]$,
    it descends to the smooth map $F\colon S^1\to S^1$.
    To see that $F$ has degree $0$ we note that it is homotopic to
    a constant map via the homotopy
    \begin{align*}
	H\colon S^1\times[0,1]&\to S^1\\
	(x,t)&\mapsto tF(x)
    \end{align*}
    (thinking of the $S^1$ in the codomain as $[0,1]$
    with its endpoints identified).
\end{soln}

\pagebreak

\begin{problem}
    [17-13]
    Let $T^2 = S^1\times S^1$ be the $2$-torus.
    Consider the two maps $f,g\colon T^2\to T^2$ given by
    $f(w,z) = (w,z)$ and $g(w,z) = (z,\ol w)$.
    Show that $f$ and $g$ have the same degree, but are not homotopic.
\end{problem}

\begin{soln}
    Since $f$ is the identity map, it has degree $1$.
    To show that $g$ has degree $1$ observe that
    $g$ is a diffeomorphism because it has a smooth inverse
    $g\inv(w,z) = (\ol z,w)$, and it is orientation-preserving because
    the determinant of the differential of $g$ at any point is
    \begin{equation*}
        \begin{vmatrix}
	    0&1\\-1&0
        \end{vmatrix} = 1,
    \end{equation*}
    and we know that orientation-preserving diffeomorphisms
    have degree $1$.
    However, these maps are not homotopic because
    $f$ induces the identity homomorphism on $\pi_1(T^2)$
    while $g$ induces the homomorphism
    \begin{equation*}
	g_*([\gamma_1],[\gamma_2])
	= ([\gamma_2], -[\gamma_1])
    \end{equation*}
    which is clearly not the identity,
    where $(\gamma_1,\gamma_2)\colon S^1\to T^2$
    is an arbitrary loop in $T^2$.
\end{soln}

\begin{problem}
    Find the de Rham cohomology groups of the $3$-torus
    $T^3 = S^1\times S^1\times S^1$.
\end{problem}

\begin{soln}
    We more generally compute $H^k(T^n)\cong\RR^{\binom nk}$
    for all nonnegative integers $n$ and integers $k$ by induction on $n$,
    starting with the observation that the singleton set $T^0$
    has de Rham cohomology
    \begin{equation*}
        H^k(T^0) = \begin{cases}
	    \RR = \RR^{\binom00} & \text{if }k = 0\\
	    0 & \text{if }k > 0.
        \end{cases}
    \end{equation*}
    Write $S^1 = U\cap V$ where $U$ and $V$ are arcs of $S^1$
    and $U\cup V = S^1$.
    Then, the intersection $U\cap V$ consists of two disjoint arcs.
    The sets $\tilde U = U\times T^{n-1}$ and $\tilde V = V\times T^{n-1}$
    give us an open cover of $T^n$, and the corresponding
    Mayer-Vietoris sequence tells us that
    \begin{equation*}
	\dots\to H^{k-1}(\tilde U\cap\tilde V)
	\overset{\delta^{k-1}}\longrightarrow H^{k}(T^n)
	\overset{\alpha^{k}}\longrightarrow
	H^{k}(\tilde U)\oplus H^{k}(\tilde V)
	\overset{\beta^{k}}\longrightarrow H^{k}(\tilde U\cap\tilde V)
	\to\dots
    \end{equation*}
    is exact for all $k$. Since $U$ and $V$ are contractible,
    $\tilde U$ and $\tilde V$ are homotopy equivalent to $T^{n-1}$.
    Also, $\tilde U\cap\tilde V$ has two components $W_1$, $W_2$
    that are homotopy equivalent to $T^{n-1}$.
    Writing $H^k(\tilde U\cap\tilde V)\cong H^k(W_1)\oplus H^k(W_2)$,
    the map $\beta^k$ is given by
    \begin{equation*}
	\eta^k([\omega],[\eta]) =
	([\omega|_{W_1}] - [\eta|_{W_1}], [\omega|_{W_2}] - [\eta|_{W_2}]),
    \end{equation*}
    so the image of this map, which equals
    the kernel of $\delta^k$, is the diagonal $([\omega],[\omega])$.
    Observe also that the image of $\alpha^k$
    (which is also the kernel of $\beta^k$)
    is also the diagonal of $H^k(T^{n-1})\oplus H^k(T^{n-1}$,
    which is isomorphic to the image of $\beta^k$.
    Thus the image and kernel of $\beta^k$ have the same dimension
    $\binom{n-1}k$, using the inductive hypothesis.

    Finally, the dimension of $H^k(T^n)$ is
    sum of the dimension of the image of $\alpha^k$
    and the dimension of the image of $\delta^{k-1}$;
    the former we know to be $\binom{n-1}k$,
    and the latter is $\dim H^{k-1}(\tilde U\cap\tilde V) =2\binom{n-1}{k-1}$
    minus $\dim\ker\delta^{k-1} = \dim\operatorname{im}\beta^{k-1}
    = \binom{n-1}{k-1}$, meaning
    \begin{equation*}
	\dim H^k(T^n) = \binom{n-1}k + \binom{n-1}{k-1} = \boxed{\binom nk}.
    \end{equation*}

    In the $n=3$ case this means
    \begin{equation*}
        H^k(T^3) = \begin{cases}
	    \RR & \text{if }k=0,3\\
	    \RR^3 & \text{if }k=1,2\\
	    0 & \text{if }k > 3.
        \end{cases}\qedhere
    \end{equation*}
\end{soln}










\end{document}
